\documentclass[14pt,a4paper,oneside]{extbook}
\usepackage[utf8]{vietnam}
\usepackage{amsmath,amssymb,mathrsfs,fancyhdr,enumerate,multirow,makecell}
\usepackage{tkz-euclide,tikz-3dplot,tikz,tkz-tab}
\usetikzlibrary{arrows.meta,calc,intersections,angles,quotes,backgrounds,shapes.geometric,patterns,through,shadows,positioning}
\usepackage{yhmath}
\usepackage[top=0.7cm, bottom=0.7cm, left=1.0cm, right=1.0cm, headheight=16pt, headsep=3mm, footskip=5mm]{geometry}
\usepackage[hidelinks,unicode]{hyperref}
\usepackage{qrcode}
\usepackage[many]{tcolorbox}
\usepackage{fontawesome}

\pagestyle{fancy}
\fancyhf{}
\fancyhead[L]{\footnotesize\bfseries\color{teal!85!black}\faBook~SỔ TAY KIẾN THỨC TOÁN 4}
\fancyhead[R]{\footnotesize\bfseries\color{blue!85!black}Lớp Toán Thầy Lĩnh -- loptoanthaylinh.vn}
\fancyfoot[C]{\footnotesize\bfseries\color{gray!80!black}-- Sổ tay kiến thức: Bài 63. Phép nhân phân số --}
\renewcommand{\headrulewidth}{0.4pt}

\begin{document}

% --- BANNER TIÊU ĐỀ BÀI HỌC ---
\noindent
\begin{tcolorbox}[
  enhanced,
  colback=teal!8!white,
  colframe=teal!80!black,
  arc=3mm, boxrule=1pt,
  left=3mm, right=3mm, top=2mm, bottom=2mm,
  halign=center
]
  {\Large\bfseries\color{teal!90!black}\faBookmark\ BÀI 63: PHÉP NHÂN PHÂN SỐ}\\[1.5pt]
  {\footnotesize\itshape\color{black!75} Quy tắc tử nhân tử, mẫu nhân mẫu, kĩ thuật rút gọn chéo trước khi nhân trên 1 trang A4}
\end{tcolorbox}

\vspace{-2.5mm}

\noindent
\begin{minipage}[t]{0.49\textwidth}\vspace{0pt}
\begin{tcolorbox}[
  enhanced,
  colback=blue!4!white,
  colframe=blue!75!black,
  arc=2mm, boxrule=0.8pt,
  left=2mm, right=2mm, top=1.8mm, bottom=1.8mm,
  title={\footnotesize\bfseries\color{white}\faTimes\ 1. Quy tắc nhân phân số}
]
{\footnotesize
Muốn nhân hai phân số, ta lấy \textbf{tử số nhân với tử số, mẫu số nhân với mẫu số}:
$$\dfrac{a}{b} \times \dfrac{c}{d} = \dfrac{a \times c}{b \times d}$$
\textbf{Nhân phân số với số tự nhiên:}\\
$$a \times \dfrac{b}{c} = \dfrac{a \times b}{c}; \quad \dfrac{a}{b} \times c = \dfrac{a \times c}{b}$$
(Lấy số tự nhiên nhân với tử số, giữ nguyên mẫu số).
}
\end{tcolorbox}
\end{minipage}
\hfill
\begin{minipage}[t]{0.49\textwidth}\vspace{0pt}
\begin{tcolorbox}[
  enhanced,
  colback=teal!4!white,
  colframe=teal!75!black,
  arc=2mm, boxrule=0.8pt,
  left=2mm, right=2mm, top=1.8mm, bottom=1.8mm,
  title={\footnotesize\bfseries\color{white}\faCut\ 2. Kĩ thuật rút gọn trước khi nhân}
]
{\footnotesize
\textbf{Bí quyết tính siêu nhanh:}
\begin{itemize}\setlength{\itemsep}{1pt}\setlength{\leftskip}{-2.5mm}
  \item Viết phép nhân dưới dạng một phân số chung $\dfrac{a \times c}{b \times d}$.
  \item Tìm các thừa số ở tử và mẫu cùng chia hết cho một số để \textbf{triệt tiêu rút gọn trước}.
  \item Nhân các số còn lại sau khi rút gọn.
  \item \textit{Lợi ích:} Kết quả nhận được luôn là \textbf{phân số tối giản}, tránh số quá lớn!
\end{itemize}
}
\end{tcolorbox}
\end{minipage}

\vspace{-2.5mm}
% --- KHỐI 3: SƠ ĐỒ TƯ DUY TIKZ ---
\noindent
\begin{tcolorbox}[
  enhanced,
  colback=orange!4!white,
  colframe=orange!80!black,
  arc=2mm, boxrule=0.8pt,
  left=2mm, right=2mm, top=1.5mm, bottom=1.5mm,
  title={\footnotesize\bfseries\color{white}\faSitemap\ SƠ ĐỒ TƯ DUY: Ý NGHĨA HÌNH HỌC PHÉP NHÂN PHÂN SỐ $\frac{1}{2} \times \frac{2}{3}$}
]
\centering
\begin{tikzpicture}[scale=0.85, transform shape]
  % Hình vuông đơn vị 1 x 1
  \draw[thick, fill=gray!8] (0,0) rectangle (6, 3);

  % Chia ngang làm 2 phần (1/2)
  \draw[dashed, thick, gray!60] (0, 1.5) -- (6, 1.5);
  \node[left, font=\scriptsize\bfseries] at (0, 0.75) {$\dfrac{1}{2}$};

  % Chia dọc làm 3 phần (2/3)
  \draw[dashed, thick, gray!60] (2, 0) -- (2, 3);
  \draw[dashed, thick, gray!60] (4, 0) -- (4, 3);
  \node[below, font=\scriptsize\bfseries] at (2, 0) {$\dfrac{2}{3}$};

  % Tô màu phần giao
  \fill[teal!40, draw=teal!80!black, thick] (0, 0) rectangle (4, 1.5);
  \node[font=\scriptsize\bfseries\color{teal!90!black}] at (2, 0.75) {Diện tích $= \dfrac{1}{2} \times \dfrac{2}{3} = \dfrac{2}{6} = \dfrac{1}{3}$};

  % Callout giải thích
  \node[draw=blue!80!black, fill=blue!8, rounded corners=2pt, font=\scriptsize\bfseries, align=left] (c1) at (9.8, 1.5) {QUY TẮC VÀNG:\\$\mathbf{\dfrac{a}{b} \times \dfrac{c}{d} = \dfrac{a \times c}{b \times d}}$\\[2pt]Tử nhân Tử\\Mẫu nhân Mẫu};
  \draw[-{Stealth}, blue!80!black, thick] (c1.west) -- (4.2, 0.75);
\end{tikzpicture}
\end{tcolorbox}

\vspace{-2.5mm}
% --- KHỐI 4: VÍ DỤ & MẸO TRÁNH BẪY ---
\noindent
\begin{minipage}[t]{0.49\textwidth}\vspace{0pt}
\begin{tcolorbox}[
  enhanced,
  colback=green!4!white,
  colframe=green!60!black,
  arc=2mm, boxrule=0.8pt,
  left=2mm, right=2mm, top=1.8mm, bottom=1.8mm,
  title={\footnotesize\bfseries\color{white}\faCheckCircle\ Ví dụ thực hành chuẩn mực}
]
{\footnotesize
\textbf{1. Nhân thông thường:}\\
$\dfrac{4}{5} \times \dfrac{3}{7} = \dfrac{4 \times 3}{5 \times 7} = \mathbf{\dfrac{12}{35}}$.\\
\textbf{2. Nhân có rút gọn chéo (siêu nhanh):}\\
$\dfrac{5}{6} \times \dfrac{3}{10} = \dfrac{5 \times 3}{6 \times 10} = \dfrac{1 \times 1}{2 \times 2} = \mathbf{\dfrac{1}{4}}$.\\
(Rút gọn $5$ với $10$ được $2$; $3$ với $6$ được $2$).\\
\textbf{3. Nhân với số tự nhiên:}\\
$3 \times \dfrac{2}{7} = \dfrac{3 \times 2}{7} = \mathbf{\dfrac{6}{7}}$.
}
\end{tcolorbox}
\end{minipage}
\hfill
\begin{minipage}[t]{0.49\textwidth}\vspace{0pt}
\begin{tcolorbox}[
  enhanced,
  colback=red!4!white,
  colframe=red!70!black,
  arc=2mm, boxrule=0.8pt,
  left=2mm, right=2mm, top=1.8mm, bottom=1.8mm,
  title={\footnotesize\bfseries\color{white}\faExclamationTriangle\ Mẹo tránh bẫy trắc nghiệm}
]
{\footnotesize
\begin{itemize}\setlength{\itemsep}{1.5pt}\setlength{\leftskip}{-2.5mm}
  \item \textbf{Bẫy quy đồng mẫu số khi nhân:}\\
  Rất nhiều học sinh máy móc đi quy đồng mẫu khi nhân phân số (SAI hoàn toàn và mất thời gian!). Phép nhân \textbf{chỉ cần lấy tử nhân tử, mẫu nhân mẫu}.
  \item \textbf{Bẫy nhân cả tử và mẫu với số tự nhiên:}\\
  $2 \times \frac{3}{5} \ne \frac{2 \times 3}{2 \times 5}$! Chỉ nhân số tự nhiên với tử số.
\end{itemize}
}
\end{tcolorbox}
\end{minipage}

\end{document}
